\documentclass[11pt,a4paper]{article}
\usepackage[utf8]{inputenc}\usepackage[T1]{fontenc}
\usepackage[provide=*,english]{babel}
\usepackage{lmodern,microtype,amsmath,amssymb,booktabs,geometry,xcolor,fancyhdr}
\usepackage[hidelinks]{hyperref}
\geometry{margin=2.3cm,headheight=14pt}
\setlength{\parindent}{0pt}\setlength{\parskip}{0.4em}
\definecolor{gcblue}{RGB}{34,73,110}\definecolor{gcgray}{RGB}{90,96,102}
\pagestyle{fancy}\fancyhf{}
\fancyhead[L]{\small\color{gcgray}GC2D: ABBA6Implicit}
\fancyhead[R]{\small\color{gcgray}Model theory}\fancyfoot[C]{\thepage}
\setlength{\emergencystretch}{3em}
\begin{document}
\begin{center}
 {\LARGE\bfseries\color{gcblue}ABBA6 Implicit}\par
 \vspace{0.35em}{\large One symmetric projection around a Yoshida composition}
\end{center}

\section{Unprojected seven-pair composition}
Let $S_{q,t}$ denote the symmetric second-order ABBA map on two spatial
copies, with signed duration $q$ and start time $t$. Set
\begin{align*}
a&=0.78451361047755726382,&b&=0.23557321335935813368,\\
c&=-1.17767998417887100695,&d&=1.31518632068391121889,
\end{align*}
and $(w_1,\ldots,w_7)=(a,b,c,d,c,b,a)$. For a trial multiplier $\mu$,
the complete base map is evaluated through
\[
 Y^{(0)}=Ez_n+N\mu,\qquad
 t_j=t_n+h\sum_{\ell<j}w_\ell,\qquad
 Y^{(j)}=S_{w_jh,t_j}(Y^{(j-1)}),\qquad
 \Psi_{h,t_n}(Y^{(0)})=Y^{(7)}.
\]
Both copies remain independent between pairs. Each pair applies an adjoint
and a direct map with weights $w_j/2$, giving fourteen stages in total.
The weights satisfy, up to round-off,
\[
 \sum_jw_j=1,\qquad \sum_jw_j^3=0,\qquad \sum_jw_j^5=0.
\]
These moments and palindromy are useful checks, but are not sufficient on
their own to prove order six: the Yoshida solution also satisfies the mixed
commutator order condition for a symmetric second-order base.

\section{One outer symmetric projection}
For $m=2N_p$ physical coordinates, with $N_p$ independent particles, define
$Ez=(z,z)$, $G=[I,-I]$, $N=G^{\mathsf T}$ and $P=\tfrac12[I,I]$.
One complete step solves
\begin{align*}
 r(\mu)&=G\Psi_{h,t_n}(Ez_n+N\mu)+2\mu=0,\\
 z_{n+1}&=P\bigl(\Psi_{h,t_n}(Ez_n+N\mu)+N\mu\bigr).
\end{align*}
The reduced formulation has $m$ unknowns. The simultaneous formulation
solves for $(Y_+,\mu)$ using
\[
 Y_+-N\mu-\Psi_{h,t_n}(Ez_n+N\mu)=0,\qquad GY_+=0.
\]
Both formulations define the same local map and permit Newton or Broyden.
There is one multiplier and one nonlinear solve per complete step. This is
exactly the ABBA4 architecture with a different unprojected recipe. It
intentionally replaces the historical composition of seven projected ABBA2
maps; previously saved trajectories and costs belong to that former map.

\section{Signed times, order and geometry}
The third and fifth pair durations are negative. Internal times can leave
$[t_n,t_n+h]$, reaching approximately $t_n-0.15759h$ and $t_n+1.15759h$;
the dynamics must be defined throughout the signed path. Adjoint stages
evaluate at their starting clock, direct stages at their ending clock.

For a smooth field, the Yoshida base has order six on doubled space. Its
exact flow preserves the diagonal, so the local diagonal defect is
$O(h^7)$. At $h=0$, $D_\mu r=4I$; the implicit function theorem therefore
gives the local multiplier $\mu=O(h^7)$. The outer projection preserves
order six. The symmetric root construction is reversible and, for the
canonical Hamiltonian splitting, symplectic on the physical diagonal.
These statements concern the local exact root. Finite stopping tolerances,
round-off and nonsmooth fields can limit the observed order and geometry.

\section{Accepted traces and passive energy}
Only the accepted root contributes to passive energy. The 28 signed shear
sources give
\[
 \Delta\kappa=\frac12\sum_j s_j[-\partial_tH(t_j,z_j)].
\]
Time and momentum do not enter the nonlinear equation. The ABBA6 observation
contains one multiplier, complete-step work metrics and seven unprojected
pair snapshots. Work histories have shape $(n_{\rm steps},1)$ rather than
$(n_{\rm steps},7)$. Differentiation composes the unprojected tangents first,
then differentiates the single outer root.

\section{Implementation correspondence and verification}
The public class is in \path{src/methods/extended/abba.py}; it inherits the
same initialization, advance and placement validation as ABBA4. Recipes and
weights live in \path{src/methods/extended/core/composition.py}. The path is
\texttt{advance}, \texttt{solve\_projection}, \texttt{compose}, then
\texttt{direct\_map} or \texttt{adjoint\_map}. Common Newton/Broyden services,
particle-block tangents and energy quadrature are shared with BM4.
Public method names remain unchanged; no forwarding modules preserve retired
internal paths. Event conversion lives in \texttt{extended/observations.py}.

Tests check the one-solve contract, fourteen signed stages, pair continuity,
nonautonomous sixth-order convergence, both residual formulations with both
solvers, reversibility, analytic tangents and symplecticity. The nonlinear
convergence test uses an independently tighter DOP853 reference. Current
architecture and diagrams are documented in
\path{docs/models/extended/simulation/extended-simulation-architecture.md}.

\clearpage
\section{Four state formulations and passive energy tracking}
For one planar particle the supported internal representations are
\[
\begin{array}{lll}
\text{physical} & z=(x,y) & \mathbb{R}^{2},\\
\text{physical with energy} & (z,t,\kappa) & \mathbb{R}^{4},\\
\text{duplicated} & (u,v),\quad u,v\in\mathbb{R}^{2} & \mathbb{R}^{4},\\
\text{duplicated with energy} & (u,v,t,\kappa) & \mathbb{R}^{6}.
\end{array}
\]
Classical methods use the physical rows; ABBA and BM4 use the duplicated rows.
\texttt{track\_energy} selects the energy variant. For $N$ planar particles the
respective dimensions are $2N$, $4N$, $4N$, and $6N$: one time entry and
one momentum per particle. All time entries equal the integration time. Classical full-cyclotron runs instead use $4N$ or $6N$ entries.
Accepted spatial copies coincide; only one physical copy is exported.
Every coordinate block contains $N$ values in component-major order. The
formulation validates particle times and aligns output round-off; the
integrator does not interpret coordinate indices.

The normalized passive balance is
\[
\dot\kappa=-\partial_tH(t,z),\qquad \kappa(t_0)=0,\qquad
\Delta K(t)=H(t,z(t))+\kappa(t)-H(t_0,z_0).
\]
Splitting accumulators from both copies are divided by two. The physical
Hamiltonian may vary with time; $\Delta K$ measures its computed balance.
Tracking leaves physical trajectories, nonlinear stopping rules and adaptive
acceptance unchanged. Dynamics must supply the momentum derivative, including
zero for an autonomous Hamiltonian.

Hairer's constraint is solely $u-v=0$: the reduced multiplier has $2N$ entries
and ABBA's simultaneous spatial solve has $6N$ unknowns. Neither $t$ nor
$\kappa$ is projected. The former eight-component projection is retired;
\texttt{state\_extension='fully\_extended'} raises migration guidance.

Dynamics, formulation, method and integration remain separate. The shared
state classes own layouts, diagonal embedding and history extraction in
\begin{quote}\small\path{src/formulations/state.py}\end{quote}
Each integration initializes a fresh method instance. Observers receive physical
states; energy histories are in \texttt{Solution.diagnostics}.
The \texttt{extended\_time} history has shape $(N,n_{\rm samples})$. Fixed-step shadow
samples leave main states and counters unchanged. Adaptive energy quadrature
uses accepted dense output; its backend state and Jacobian remain physical.


\end{document}
